\chapter{Algorithm 1: Iteratively Reweighted Least Squares}
\label{chap:irls}

Iteratively Reweighted Least Squares (IRLS) \cite{daubechies2008iterativelyreweightedsquaresminimization} is a classical algorithm for solving
optimization problems involving non-smooth objective functions of the form of a $p$-norm.
It handles non-smoothness
by solving a sequence of \textit{weighted} least-squares problems, each of which is
smooth and admits a closed-form solution.

% ------------------------------------------------------------------
\section{A strict upper bound for the \texorpdfstring{$L_1$}{L1} norm}
\label{sec:irls_idea}

We solve the differentiability problem replacing the $L_1$ term with a smooth surrogate function, using the inequality $(|w_i| - |(w_k)_i|)^2 \ge 0$, which, after expansion and division by $2|(w_k)_i|$, gives the majorization:
\begin{equation}
    |w_i| \le \frac{w_i^2}{2|(w_k)_i|} + \frac{|(w_k)_i|}{2} \qquad (w_k)_i \neq 0
    \label{eq:irls_majorization_inequality}
\end{equation}
Summing over $i = 1, \ldots, H$ returns the bound on the $L_1$ norm
\begin{equation}
    \|\mathbf{w}\|_1 \le \frac{1}{2} \mathbf{w}^T \mathbf{W}_k^T\mathbf{W}_k \mathbf{w} + \frac{1}{2}\|\mathbf{w}_k\|_1
    \label{eq:irls_quadratic}
\end{equation}
where $\mathbf{W}_k \in \mathbb{R}^{H \times H}$ is a diagonal matrix with entries:
\begin{equation}
    (\mathbf{W}_k)_{ii} = |(w_{k})_i|^{-1/2}
    \label{eq:irls_weights_exact}
\end{equation}
This bounds the non-smooth $L_1$ term using a \textit{smooth, quadratic} expression
$\frac{1}{2}\mathbf{w}^T \mathbf{W}_k^T \mathbf{W}_k \mathbf{w}$, plus a constant term that depends only on the current iteration $\mathbf{w}_k$.
$\mathbf{W}_k^T \mathbf{W}_k$ assigns each component of $\mathbf{w}$ a per-iteration weight $|(w_k)_i|^{-1}$ \cite{chartrand2008iteratively}: small components are pushed further toward zero by a larger weight, large ones are left almost free.

\paragraph{Handling near-zero components.}
Equation~\eqref{eq:irls_weights_exact} becomes numerically unstable when $(w_k)_i
\approx 0$, since $|(w_k)_i|^{-1/2}$ can blow up. To prevent division by (near) zero,
a \textbf{safety threshold} $\varepsilon_{\mathrm{thr}} > 0$ is introduced:
\begin{equation}
    (\mathbf{W}_k)_{ii} = \Bigl(\max\!\bigl(|(w_k)_i|,\, \varepsilon_{\mathrm{thr}}\bigr)\Bigr)^{-1/2}
    \label{eq:irls_weights}
\end{equation}
Typical values are $\varepsilon_{\mathrm{thr}} \in [10^{-6}, 10^{-10}]$.

% ------------------------------------------------------------------
\section{The Majorization-Minimization interpretation}
\label{sec:irls_mm}

IRLS fits inside the Majorization-Minimization (MM) framework \cite{nguyen2016introductionmmalgorithmsmachine}. Instead of minimizing $f$ directly, at each iteration $k$ we build a \emph{surrogate} $Q(\mathbf{w}, \mathbf{w}_k)$ function that:
\begin{enumerate}
    \item Majorizes $f$: $Q(\mathbf{w}, \mathbf{w}_k) \ge f(\mathbf{w})$ for all $\mathbf{w}$
    \item Touches $f$ at $\mathbf{w}_k$: $Q(\mathbf{w}_k, \mathbf{w}_k) = f(\mathbf{w}_k)$
    \item Is smooth, so that it can be minimized in closed form
\end{enumerate}
Plugging the quadratic upper bound~\eqref{eq:irls_quadratic} into~\eqref{eq:obj_fun} gives the following surrogate:
\begin{equation}
    Q(\mathbf{w}, \mathbf{w}_k) = \frac{1}{2}\|\mathbf{Xw} - \mathbf{y}\|_2^2
    + \frac{\lambda_{\text{LASSO}}}{2}\|\mathbf{W}_k\mathbf{w}\|_2^2
    + \frac{\lambda_{\text{LASSO}}}{2}\|\mathbf{w}_k\|_1
    \label{eq:irls_surrogate}
\end{equation}
We need to check that the properties are satisfied:
\begin{enumerate}
    \item Majorization property: follows from \eqref{eq:irls_majorization_inequality}
    \item Touching property: at $\mathbf{w} = \mathbf{w}_k$, we get the following $Q$ function
        \[
            Q(\mathbf{w}_k, \mathbf{w}_k)
            = \tfrac{1}{2}\|\mathbf{Xw}_k - \mathbf{y}\|_2^2
            + \tfrac{\lambda_{\text{LASSO}}}{2}\|\mathbf{w}_k\|_1
            + \tfrac{\lambda_{\text{LASSO}}}{2}\|\mathbf{w}_k\|_1
            = f(\mathbf{w}_k)
        \]
    \item \textbf{Smoothness property:} The surrogate $Q(\mathbf{w}, \mathbf{w}_k)$ is a function of $\mathbf{w}$, since $\mathbf{w}_k$ is fixed (constant). Its gradient is 
    \[\nabla_{\mathbf{w}} Q(\mathbf{w}, \mathbf{w}_k) = \left(\mathbf{X}^T\mathbf{X} + \lambda_{\text{LASSO}}\mathbf{W}_k^T\mathbf{W}_k\right)\mathbf{w} - \mathbf{X}^T\mathbf{y}\]
    Since it is an affine function of $\mathbf{w}$, the Hessian is a constant matrix. Because of that, $Q$ is infinitely differentiable ($C^\infty$), and also globally $L$-smooth (with a Lipschitz continuous gradient).
\end{enumerate}

Taking $\mathbf{w}_{k+1} = \arg\min_{\mathbf{w}} Q(\mathbf{w}, \mathbf{w}_k)$, the first two MM properties combined give $f(\mathbf{w}_{k+1}) \le Q(\mathbf{w}_{k+1}, \mathbf{w}_k) \le Q(\mathbf{w}_k, \mathbf{w}_k) = f(\mathbf{w}_k)$, so the sequence $\{f(\mathbf{w}_k)\}$ is non-increasing.

% ------------------------------------------------------------------
\section{The weighted least-squares subproblem}
\label{sec:irls_subproblem}

To find $\mathbf{w}_{k+1}$, we must minimize the surrogate function $Q(\mathbf{w}, \mathbf{w}_k)$ with respect to $\mathbf{w}$:
$$\mathbf{w}_{k+1} = \underset{\mathbf{w}}{\arg\min} \left( \frac{1}{2}\|\mathbf{Xw} - \mathbf{y}\|_2^2 + \frac{\lambda_{\text{LASSO}}}{2}\|\mathbf{W}_k\mathbf{w}\|_2^2 + \frac{\lambda_{\text{LASSO}}}{2}\|\mathbf{w}_k\|_1 \right)$$
Since the term $\frac{\lambda_{\text{LASSO}}}{2}\|\mathbf{w}_k\|_1$ is constant with respect to $\mathbf{w}$, it does not affect the optimization and can be dropped. The minimization problem simplifies to a standard \textbf{weighted least-squares problem with $L_2$ regularization}:
\begin{equation}
    \mathbf{w}_{k+1} = \underset{\mathbf{w}}{\arg\min} \left( \frac{1}{2}\|\mathbf{Xw} - \mathbf{y}\|_2^2 + \lambda_{\text{IRLS}} \|\mathbf{W}_k \mathbf{w}\|^2_2 \right)
    \label{eq:irls_subproblem}
\end{equation}
using:
\begin{equation}
    \lambda_{\text{IRLS}} = \frac{1}{2}\lambda_{\text{LASSO}}
    \label{eq:lambda_relation}
\end{equation}

% \paragraph{Relation to the LASSO regularization parameter.}\mbox{} \\
% The factor $1/2$ in~\eqref{eq:lambda_relation} is forced by the majorization constant in~\eqref{eq:irls_quadratic}. The upper bound on $\|\mathbf{w}\|_1$ carries a coefficient $\tfrac{1}{2}$ in front of the quadratic form $\mathbf{w}^T\mathbf{W}_k^T\mathbf{W}_k\mathbf{w}$, so the penalty entering the surrogate~\eqref{eq:irls_surrogate} is $\tfrac{\lambda_{\text{LASSO}}}{2}\|\mathbf{W}_k\mathbf{w}\|_2^2$. Rewriting this as $\lambda_{\text{IRLS}}\|\mathbf{W}_k\mathbf{w}\|_2^2$ in~\eqref{eq:irls_subproblem} therefore requires $\lambda_{\text{IRLS}} = \tfrac{1}{2}\lambda_{\text{LASSO}}$. This is precisely the factor that makes the fixed point of IRLS coincide with the LASSO optimum, as shown in the proof of Theorem~\ref{thm:irls_convergence}.

\begin{proposition}[Weighted normal equation solution]
    \label{prop:irls_normal}
    The unique solution of~\eqref{eq:irls_subproblem} satisfies:
    \begin{equation}
        \bigl(\mathbf{X}^T\mathbf{X} + 2\lambda_{\text{IRLS}}\mathbf{W}_k^T \mathbf{W}_k \bigr)\,
        \mathbf{w}_{k+1} = \mathbf{X}^T\mathbf{y}
        \label{eq:irls_normal}
    \end{equation}
\end{proposition}

\begin{proof}
    The gradient of the objective function \eqref{eq:irls_subproblem} is:
    $$\nabla f_k(\mathbf{w}) = \mathbf{X}^T(\mathbf{Xw} - \mathbf{y}) + 2\lambda_{\text{IRLS}}\mathbf{W}_k^T \mathbf{W}_k\mathbf{w}$$
    Setting this to zero results in \eqref{eq:irls_normal}.
    The coefficient matrix $\mathbf{X}^T\mathbf{X} + 2\lambda_{\text{IRLS}}\mathbf{W}_k^T\mathbf{W}_k$
    is symmetric positive definite (it is the sum of $\mathbf{X}^T\mathbf{X} \succeq 0$
    and $2\lambda_{\text{IRLS}}\mathbf{W}_k^T\mathbf{W}_k \succ 0$ since $\varepsilon_{\mathrm{thr}} > 0$),
    so the system has a unique solution.
\end{proof}

Since $\mathbf{W}_k$ is diagonal, defining $\mathbf{Q}_k = \mathbf{X}^T\mathbf{X}
+2\lambda_{\text{IRLS}}\mathbf{W}_k^T\mathbf{W}_k$ and $\mathbf{b} = \mathbf{X}^T\mathbf{y}$,
equation~\eqref{eq:irls_normal} reduces to:
\begin{equation}
    \mathbf{Q}_k\, \mathbf{w}_{k+1} = \mathbf{b}
    \label{eq:irls_linear_system}
\end{equation}
This is an $H \times H$ symmetric positive definite linear system, which can
be solved efficiently at each iteration. Note that $\mathbf{A} = \mathbf{X}^T\mathbf{X}$ and $\mathbf{b} = \mathbf{X}^T\mathbf{y}$
can be pre-computed once before the loop, so only the diagonal $2\lambda_{\text{IRLS}}\mathbf{W}_k^T\mathbf{W}_k$ is updated at each iteration.

\paragraph{Solving the linear system.}
Because $\mathbf{Q}_k$ is SPD, two efficient methods can be used:
\begin{itemize}
    \item \textbf{Cholesky factorization:}~\cite{trefethen1997numerical} computes $\mathbf{Q}_k = \mathbf{L}\mathbf{L}^T$
    in $O(H^3/3)$ operations, then solves two triangular systems in $O(H^2)$. This
    is the direct method of choice for moderate-sized problems.
    \item \textbf{Conjugate Gradient (CG):}~\cite{trefethen1997numerical} an iterative method that converges in at
    most $H$ steps in exact arithmetic, with convergence rate depending on the condition
    number $\kappa(\mathbf{Q}_k)$. Preferable for large-scale problems where $H$ is
    too large for Cholesky.
\end{itemize}

% ------------------------------------------------------------------
\section{The complete algorithm}
\label{sec:irls_algorithm}

\begin{algorithm}[H]
    \caption{Iteratively Reweighted Least Squares (IRLS)}
    \label{alg:irls}
    \begin{algorithmic}[1]
        \Require $\mathbf{X} \in \mathbb{R}^{M \times H}$, $\mathbf{y} \in \mathbb{R}^M$,
        $\lambda_{\text{LASSO}} > 0$, $\varepsilon_{\mathrm{thr}} > 0$, $\varepsilon_{\mathrm{stop}} > 0$,
        $k_{\max}$
        \State Pre-compute $\mathbf{A} \leftarrow \mathbf{X}^T\mathbf{X}$,
        $\mathbf{b} \leftarrow \mathbf{X}^T\mathbf{y}$
        \State Define $\lambda_{\text{IRLS}} \leftarrow \frac{1}{2}\lambda_{\text{LASSO}}$
        \State Initialize $\mathbf{w}_0 \leftarrow \mathbf{A}^{-1}\mathbf{b}$
        (ordinary least-squares solution)
        \For{$k = 0, 1, 2, \ldots, k_{\max}$}
            \State Compute weights: $(\mathbf{W}_k)_{ii} \leftarrow
            \bigl(\max(|(w_k)_i|,\, \varepsilon_{\mathrm{thr}})\bigr)^{-1/2}$ for all $i$
            \State Assemble system: $\mathbf{Q}_k \leftarrow \mathbf{A}
            + 2\lambda_{\text{IRLS}}\mathbf{W}_k^T\mathbf{W}_k$
            \State Solve: $\mathbf{Q}_k\,\mathbf{w}_{k+1} = \mathbf{b}$
            \quad \textit{(via Cholesky or CG)}
            \If{$\|\mathbf{w}_{k+1} - \mathbf{w}_k\| \,/\, \max(1, \|\mathbf{w}_k\|)
            < \varepsilon_{\mathrm{stop}}$}
                \State \textbf{break} \quad \textit{(convergence reached)}
            \EndIf
        \EndFor
        \State \Return $\mathbf{w}_{k+1}$
    \end{algorithmic}
\end{algorithm}

\paragraph{Initialization.} We initialize $\mathbf{w}_0$ with the unregularized least-squares solution $\mathbf{A}^{-1}\mathbf{b}$ (\textbf{warm start}): the weights $(\mathbf{W}_0)_{ii} = |(w_0)_i|^{-1/2}$ are well defined without having to lean on $\varepsilon_{\mathrm{thr}}$ at the first step. A \textbf{cold start} with $\mathbf{w}_0 = \mathbf{0}$ would instead force every weight onto the threshold and slow the first iterations down. In the implementation we add a vanishing ridge $10^{-12}\mathbf{I}$ to $\mathbf{A}$ before this first solve: numerically indistinguishable from $\mathbf{A}^{-1}\mathbf{b}$ when $\mathbf{A} \succ 0$, it keeps the warm start well-posed if $\mathbf{X}$ is rank-deficient.

\paragraph{Stopping criterion.} The algorithm ends when the relative change between successive iterations goes below $\varepsilon_{\mathrm{stop}}$:
$$\frac{\|\mathbf{w}_{k+1} - \mathbf{w}_k\|}{\max(1, \|\mathbf{w}_k\|)} < \varepsilon_{\mathrm{stop}}$$
where the $\max(1, \|\mathbf{w}_k\|)$ normalization makes the test scale-invariant and prevents false triggers when $\|\mathbf{w}_k\|$ is small.

% ------------------------------------------------------------------
\section{Convergence analysis}
\label{sec:irls_convergence}

\paragraph{Monotone decrease.}
IRLS monotonically decreases the function. Since $f$ is bounded below by zero, the sequence $\{f(\mathbf{w}_k)\}$ converges.

\begin{theorem}[Fixed-point consistency]
    \label{thm:irls_convergence}
    If the sequence $\{\mathbf{w}_k\}$ generated by IRLS converges to a point
    $\mathbf{w}^*$ such that $|(w^*)_i| > \varepsilon_{\mathrm{thr}}$ for all $i$,
    then $\mathbf{w}^*$ satisfies~\eqref{eq:optimality}.
\end{theorem}

\begin{proof}
At a fixed point $\mathbf{w}_k = \mathbf{w}_{k+1} = \mathbf{w}^*$, \eqref{eq:irls_normal} gives:
\[
    \mathbf{X}^T(\mathbf{X}\mathbf{w}^* - \mathbf{y})
    + 2\lambda_{\text{IRLS}}\,\mathbf{W}_*^T\mathbf{W}_*\mathbf{w}^* = \mathbf{0}
\]
Since $|(w^*)_i| > \varepsilon_{\text{thr}}$ for all $i$, the threshold is inactive and
$(\mathbf{W}_*^T\mathbf{W}_*)_{ii} = |(w^*)_i|^{-1}$.
The $i$-th component of the regularizer term is:
\[
    (\mathbf{W}_*^T\mathbf{W}_*\mathbf{w}^*)_i
    = \frac{(w^*)_i}{|(w^*)_i|}
    = \mathrm{sign}((w^*)_i)
\]
Using \eqref{eq:lambda_relation}:
\[
    \mathbf{X}^T(\mathbf{X}\mathbf{w}^* - \mathbf{y})
    + \lambda_{\text{LASSO}}\,\mathrm{sign}(\mathbf{w}^*) = \mathbf{0}
\]
Since $|(w^*)_i| \neq 0$ for all $i$, the $L_1$ norm is differentiable at
$\mathbf{w}^*$, so $\partial\|\mathbf{w}^*\|_1 = \mathrm{sign}(\mathbf{w}^*)$
\cite{frangioni-slides-nonsmooth}. By the sum rule for subdifferentials
\cite{frangioni-slides-unconstrained}, the equation above is exactly
$\mathbf{0} \in \partial f(\mathbf{w}^*)$, the optimality condition
in \eqref{eq:optimality} is met.
\end{proof}


\paragraph{Convergence rate.}
We expect a linear rate empirically: $\|\mathbf{w}_{k+1} - \mathbf{w}^*\|$ contracts by a factor governed by the conditioning of $\mathbf{Q}_k = \mathbf{X}^T\mathbf{X} + 2\lambda_{\text{IRLS}}\mathbf{W}_k^T\mathbf{W}_k$, whose second term regularises the small components and improves the conditioning of $\mathbf{X}^T\mathbf{X}$ via a diagonal shift. A formal linear-rate result for IRLS on $\ell_q$ minimisation ($q \le 1$) appears in~\cite[Theorem~5.3]{daubechies2008iterativelyreweightedsquaresminimization} under a \emph{null space property} (NSP) on the matrix~$\mathbf{X}$ (no sparse vector is in the kernel of $\mathbf{X}$). That hypothesis does not hold in general for ELM features (no sparsity-recovery guarantee is claimed for the random $\mathbf{W}_1$ regime), so we use Daubechies et al.\ only as the closest available analysis. The linear rate is verified empirically in Section~\ref{sec:convergence}.

\paragraph{Verification of hypotheses for the ELM problem.}
The two assumptions of Theorem~\ref{thm:irls_convergence} are satisfied for~\eqref{eq:obj_fun} as follows.
\begin{enumerate}
\item \emph{Convergence of the sequence.} Under the assumption that $\mathbf{X}$ has full column rank, $f$ is strictly convex and admits a unique minimizer $\mathbf{w}^*$ (Proposition~\ref{prop:irls_prop}). IRLS monotonically decreases $f$ (Section~\ref{sec:irls_mm}) and the level sets of $f$ are compact by coercivity, so $\{\mathbf{w}_k\}$ has at least one accumulation point. The MM descent identity forces any accumulation point $\bar{\mathbf{w}}$ to be a fixed point of the IRLS map. By Theorem~\ref{thm:irls_convergence} every fixed point with $|(\bar{w})_i| > \varepsilon_{\mathrm{thr}}$ for all $i$ satisfies the optimality condition, and by uniqueness $\bar{\mathbf{w}} = \mathbf{w}^*$, so the whole sequence converges.
\item \emph{Non-zero components at convergence.} The hypothesis $|(w^*)_i| > \varepsilon_{\mathrm{thr}}$ fails exactly on the inactive set $I = \{i : (w^*)_i = 0\}$ of the sparse LASSO solution. Strictly, IRLS therefore does not solve the original problem~\eqref{eq:obj_fun} but the smoothed variant induced by the clip~\eqref{eq:irls_weights} on the weights, in which the $L_1$ term is replaced by a surrogate that agrees with $|w_i|$ for $|w_i| > \varepsilon_{\mathrm{thr}}$ and with a quadratic otherwise. The minimiser of this smoothed problem differs from the LASSO minimiser by at most $O(\varepsilon_{\mathrm{thr}})$ in objective and recovers the LASSO solution as $\varepsilon_{\mathrm{thr}} \to 0$. On the active set $A = \{i : |(w^*)_i| > \varepsilon_{\mathrm{thr}}\}$ Theorem~\ref{thm:irls_convergence} applies and reproduces the LASSO optimality condition exactly. On $I$ the smoothed gradient at $(\bar{w})_i \le \varepsilon_{\mathrm{thr}}$ pins the component near zero and the corresponding row of~\eqref{eq:optimality} is satisfied by some $s_i \in [-1,1]$, in agreement with the subdifferential of $|w_i|$ at $0$. The bound $\varepsilon_{\mathrm{thr}} = 10^{-8}$ that we adopt (Section~\ref{sec:params}) is tight enough that this approximation is not detectable against the Clarabel reference (Section~\ref{sec:real_data}).
\end{enumerate}

% ------------------------------------------------------------------
\section{Computational cost}
\label{sec:irls_cost}

The computational cost of each iteration of IRLS is the following:
\begin{itemize}
    \item \textbf{Weight update} $\mathbf{W}_k$: $O(H)$ simply iterating over
    the components of $\mathbf{w}_k$.
    \item \textbf{Matrix update} $\mathbf{Q}_k = \mathbf{A} + 2\lambda_{\text{IRLS}}\mathbf{W}_k^T\mathbf{W}_k$: $O(H)$, since $\mathbf{A} = \mathbf{X}^T\mathbf{X}$ is pre-computed and only
    the diagonal changes.
    \item \textbf{Linear system solve} $\mathbf{Q}_k \mathbf{w}_{k+1} = \mathbf{b}$:
    $O(H^3/3)$ with Cholesky, or $O(pH^2)$ with CG for $p$ CG steps (in the following we assume using Cholesky).
\end{itemize}
Also, we need $O(MH^2)$ for the $\mathbf{A}$ matrix computation and $O(MH)$ for the $\mathbf{b}$ vector creation, and, in case of warm start, we need the OLS solution $\mathbf{w}_0 \leftarrow \mathbf{A}^{-1}\mathbf{b}$, computed in $O(H^3)$.
The dominant cost per iteration is the linear system solve: $O(H^3)$ with Cholesky.
However, the total number of IRLS iterations is typically small (50--100), and the
pre-computation costs get amortized over all iterations. The overall cost is therefore:
$$\underbrace{O(MH^2) + O(MH) + O(H^3)}_{\text{pre-computation}} + \underbrace{k_{\mathrm{IRLS}} \cdot O(H^3)}_{\text{iterations cost}}$$
For problems where $H \ll M$, the pre-computation dominates and IRLS is very efficient.

A cold start $\mathbf{w}_0 = \mathbf{0}$ saves the $O(H^3)$ OLS solve in the pre-computation but loses iterations to compensate, since the initial distance $\|\mathbf{w}_0 - \mathbf{w}^*\|$ multiplies the rate of Theorem~\ref{thm:irls_convergence}. We use the warm start by default and quantify the trade-off in Section~\ref{sec:warm_cold}.

% ------------------------------------------------------------------
\section{Expected behavior on our problem}
\label{sec:irls_expected}

On the ELM LASSO problem we expect $f(\mathbf{w}_k)$ to decrease at every iteration. Any non-decreasing step indicates an implementation error. On a semilog plot of $f(\mathbf{w}_k) - f^*$ against $k$ the curve should be approximately straight, with a slope governed by the conditioning of $\mathbf{Q}_k$ and the safety threshold $\varepsilon_{\mathrm{thr}}$.

Sparsity arises dynamically: components $(w_k)_i$ that start small receive larger weights $|(w_k)_i|^{-1}$ and are pushed further toward zero. At convergence many components lie inside the $\varepsilon_{\mathrm{thr}}$-ball of zero. The algorithm has two hyperparameters: $\lambda_{\text{LASSO}}$ (model) and $\varepsilon_{\mathrm{thr}}$ (numerical). Larger $\lambda_{\text{LASSO}}$ gives sparser solutions and a higher training residual. It also enters $\mathbf{Q}_k$'s conditioning and so affects the rate. Smaller $\varepsilon_{\mathrm{thr}}$ approximates the $L_1$ term more tightly at the cost of stability when a component approaches zero. Default: $\varepsilon_{\mathrm{thr}} = 10^{-8}$ (calibrated in Section~\ref{sec:params}).